\documentclass[10pt]{article}
\usepackage[margin=0.78in]{geometry}
\usepackage[T1]{fontenc}
\usepackage{amsmath,amssymb,amsthm,booktabs,hyperref}
\hypersetup{colorlinks=true,urlcolor=blue,linkcolor=blue}
\title{\textbf{W33 Walsh-Qutrit Hardware Phase II}\\
Exact $PU(3)$ Density, a Four-Rail PIC Test Cell,\\
Calibrated Phase Noise, 160-Incidence Transfer,\\
UP5K Placement and Two-Photon Entanglement}
\author{W33-Theory photonic architecture programme}
\date{23 October 2026 --- software and layout reference, not foundry signoff}
\newtheorem{theorem}{Theorem}
\begin{document}\maketitle
\begin{abstract}
We implement five executable studies that connect the $W(3,3)$ four-port
line-bus architecture to photonic quantum information. Their strongest
mathematical finding is an exact proof that the qutrit Fourier gate
$F_3$ and the line-bus Walsh sign reflection $D=\mathrm{diag}(-1,1,-1)$
generate a dense subgroup of the projective single-qutrit group $PU(3)$.
The proof uses an infinite-order algebraic trace witness and an exact
modulo-13 certificate that the joint adjoint commutant has dimension one.
On the engineering side we generate a process-neutral four-rail GDS
floorplan; a detector-calibrated synthetic phase-noise experiment;
160 distinct W33 point--line incidence mode transfers; and run
real YoWASP Yosys/nextpnr iCE40UP5K place and route of the dual
optical/logical frame controller.
Finally, a two-photon, eight-rail passive coupler with coincidence
and logical-code postselection produces a mathematically entangled
pair of qutrit states. None of these software results constitutes a
fabricated optical device or a composable fault-tolerant quantum computer.
\end{abstract}
\tableofcontents

\section{The inherited Walsh-qutrit architecture}
The W33 graph has 40 points, 40 lines of four points each, and
160 point--line incidences. Each four-port optical line bus has
the ideal XOR permutations $P_m$, $m\in\mathbb F_2^2$.
Let $W=H_2\otimes H_2$, let $a$ be its uniform
column and $V$ its other three columns. Then
\[
V^\dagger V=I_3,\quad VV^\dagger+aa^\dagger=I_4,\quad
P_mV=VD_m,\quad
D_m=\mathrm{diag}_{s\in\{1,2,3\}}(-1)^{s\cdot m}.
\]
This is a four-rail coherent \emph{one-photon qutrit}, not
the naive three-physical-port encoding; its ideal auxiliary
leakage under XOR masks is zero. The actual MZI phase frame
also contributes $i^{w(m)}$ and is tracked separately in
the four-bit $\mathbb Z_4\times\mathbb F_2^2$ controller.

\section{Track 5: a rigorous ideal single-qutrit density theorem}
Write $\omega=e^{2\pi i/3}$ and
\[
F=\frac1{\sqrt3}(\omega^{jk})_{j,k=0}^2,\qquad
D=\mathrm{diag}(-1,1,-1),\qquad A=FD.
\]
\begin{theorem}
The subgroup of $PU(3)$ generated by the projective classes
of $F$ and $D$ is dense in $PU(3)$.
\end{theorem}
\begin{proof}
Exact symbolic calculation yields
$\operatorname{tr}(A)=-1/\sqrt3$ and $\det A=-i$.
The trace has norm $-1/3$ in $\mathbb Q(\sqrt3)$, so it
is not an algebraic integer. Finite projective order would
force each eigenvalue of $A$ to be a root of unity
(because its determinant is a root of unity), forcing
the trace to be an algebraic integer. Thus $A$ has
infinite projective order.

Let $K$ be the closure of $\langle F,D\rangle$ in $PU(3)$.
Its identity-component Lie algebra $\mathfrak k$ is nonzero,
because a compact zero-dimensional Lie group is finite.
As a normal identity component, $\mathfrak k$ is invariant
under the adjoint actions of $F$ and $D$ on
$\mathfrak{su}(3)$, an eight-dimensional real Euclidean space.

Take the eight normalized Gell--Mann Hermitian matrices as
basis and compute the two $8\times8$ real orthogonal adjoint
matrices. Stack the commutant constraints
\[
(I_8\otimes\mathrm{Ad}(U)
 -\mathrm{Ad}(U)^{\mathsf T}\otimes I_8)
 \operatorname{vec}(T)=0,\qquad U\in\{F,D\},
\]
into a $128\times64$ system $C$. Every entry belongs to
$\mathbb Q(\sqrt3)$ with denominators not divisible by 13.
Reducing $\sqrt3\mapsto4$ in $\mathbb F_{13}$, since
$4^2\equiv3\pmod{13}$, gives
$\operatorname{rank}_{\mathbb F_{13}}C=63$
with pivot columns $0,\ldots,62$. This certifies
$\operatorname{rank}_{\mathbb R}C\ge63$.
The identity operator always commutes with both adjoints,
so the characteristic-zero rank is exactly 63 and the
joint real commutant is precisely the scalar matrices.

A proper nonzero real invariant subspace for an orthogonal
representation would have an orthogonal projection
commuting with both adjoint operators, contradicting
the scalar commutant. Hence the adjoint action is
real-irreducible. Since $\mathfrak k$ is nonzero and
adjoint-invariant, it equals $\mathfrak{su}(3)$.
Therefore the compact closure $K$ is all of $PU(3)$.
\end{proof}
This is an independent exact certificate for the
specific generator pair, \emph{not} a claim that
approximate Clifford+$R$ universality was unknown:
qutrit metaplectic gate synthesis is established prior art.
The theorem does not supply gate counts for a prescribed
precision, a multi-qutrit entangler, hardware fidelity,
or physical realization of either ideal generator.

\section{Track 1: four-rail PIC validation cell}
The file
\texttt{hardware/generated/w33\_walsh\_qutrit\_lab\_ref\_20261023.gds}
has a four-rail waveguide reference skeleton,
nominal 20-$\mu$m pitch, 0.5-$\mu$m guide width and
850-$\mu$m body length. In symbolic keep-out layer 99,
it labels a Walsh encoder, XOR bus, additional physical
$S_4$ rail-three-cycle, encoded $F_3$ tile and Walsh decoder.
After decoding, output rail 0 is the uniform auxiliary
character for leakage monitoring; rails 1--3 are the
logical qutrit outputs. The mathematical circuit
\[
U_{\rm target}=W[W\operatorname{diag}(1,F_3)W]\,
\Pi_{\rm 3cycle}\,P_m\,W
\]
is verified unitary and auxiliary-preserving for four
masks and either three-cycle state. The process-neutral
GDS \emph{does not contain physically functioning}
directional couplers, modulators or crossed waveguides.
It has no foundry PDK, DRC, LVS or actual circuit-level
complex $S$-parameter simulation. This is the required
annotated geometry handoff, not a tapeout deliverable.

\section{Track 2: correlated phase noise and falsifiable counts}
For a maximally mixed encoded qutrit, differential rail
phase noise obeys
\[
p_{\mathrm{aux}}=
\frac13\left(1-\left|
\frac14\sum_{j=0}^{3}e^{i\phi_j}\right|^2\right).
\]
If the four phases are zero-mean Gaussian, have variance
$\sigma^2$ and uniform correlation coefficient
$0\le\rho\le1$, then
\[
\boxed{\mathbb E[p_{\mathrm{aux}}]
  =\frac{1-e^{-\sigma^2(1-\rho)}}4.}
\]
Thus correlated common-mode fluctuations cancel exactly.
Mean auxiliary counts alone identify only
$\sigma_{\mathrm{diff}}^2=\sigma^2(1-\rho)$;
a separate differential phase monitor is required to
measure $\sigma$ and $\rho$ independently.

With detector efficiency $\eta$ and per-trial dark
false-click probability $d$, a simplified click model is
$p_{\rm click}=d+(1-d)\eta p_{\rm aux}$.
The reproducible software generates synthetic
phase-realization and binomial detector-count tables,
and supplies an explicitly separate CSV interface
for actual future lab counts. It does not invent
laboratory measurements; detector saturation,
waveguide loss, indistinguishability and calibration
uncertainty must be independently characterized.

\section{Track 3: 160-incidence inter-bus state transfer}
The graph-level 40-mode overlap model and the physical
point--line incidence architecture are \emph{not}
the same Hilbert-space identification. For the
160 independent incidence rails $(L,p)$, even two
lines sharing W33 point $p$ possess two distinct
physical rail labels $(L,p)$ and $(L',p)$.
The exact ideal transfer of the complete four-rail
Walsh carrier between any two lines is a permutation
of \emph{four distinct rail pairs}.
The script enumerates all $\binom{40}{2}=780$
line pairs, recovering 240 intersecting and
540 disjoint pairs, and verifies zero residual
on the three-dimensional encoded isometry.
The 160-incidence adjacency model has
240 potential intra-bus links and 240 possible
inter-bus point-junction links, each incidence
degree six. These are \emph{logical possible
connections}; junction optical couplers, waveguides,
clock synchronization and collision-free four-rail
paths are not yet implemented.

\section{Track 4: synthesizable and placed dual frame}
The existing \texttt{w33\_walsh\_qutrit\_frame.sv}
maintains two bits of scalar phase in $\mathbb Z_4$
and two bits of logical Walsh mask in $\mathbb F_2^2$.
Icarus Verilog ran 1200 randomized cycles, with
368 accepted commands and 416 stalls, without
a scoreboard mismatch. YoWASP Yosys produced a
nine-\texttt{SB\_LUT4}, eight-\texttt{SB\_DFFER}
technology netlist. YoWASP nextpnr placed and
routed it for a \texttt{UP5K-sg48} reference
with 15/5280 logic cells and 19/39 IO,
one global buffer and a 24 MHz timing target.
Its reported maximum clock frequency was
\textbf{78.91 MHz}, with a PASS at 24 MHz.
The tool \emph{automatically placed every IO};
no actual board PCF, electrodes, waveguide
phase controllers, time-tagged detectors,
photonic source or external timing was connected.
This is a legitimate unconstrained reference
place-and-route result, \emph{not} a deployed FPGA
or a measured optical-computation rate.

\section{Track 5 extension: an eight-rail, two-photon primitive}
The ideal single-qutrit density theorem is
insufficient for multi-qutrit computation.
The abstract nine-dimensional
controlled phase $\operatorname{diag}(\omega^{jk})$
entangles $|+\rangle_3\otimes|+\rangle_3$
maximally, but is not implemented by
independent single-photon local tiles.

A more concrete photonic target consists of
\emph{two indistinguishable photons} in two
independent four-rail Walsh codes. Add a
directional coupler only between their
respective physical rail-0 modes; let its
single-photon amplitude mixer be
\[
\begin{pmatrix} \cos\theta&i\sin\theta\\
i\sin\theta&\cos\theta\end{pmatrix}.
\]
For input modes $j$ in bank A and $k$ in bank B,
the one-photon-per-bank scattering amplitude is
a two-boson permanent,
\[
K_{ab,jk}=U_{a,j}U_{4+b,4+k}
         +U_{a,4+k}U_{4+b,j}.
\]
Project first into one photon per four-rail bank
and then both three-dimensional Walsh
logical code subspaces. For a product
logical input $|0\rangle_3\otimes|0\rangle_3$,
the ideal full-exchange setting
$\theta=\pi/2$ yields:
\[
p_{\text{one photon each bank}}=\frac58,\quad
p_{\text{both in code}}=\frac12,\quad
\mathcal N_{\text{conditional}}=\frac14,
\]
where $\mathcal N$ is partial-transpose negativity.
The numerical calculation was checked using
the complete eight-mode two-boson transfer matrix
and independent tests. This constructs a
\emph{postselected entangled-state preparation
target}, not a deterministic or composable
two-qutrit gate: ordinary photon counting
destroys the encoded photons. Any useful
online quantum gate requires a genuine
heralding/QND or measurement-induced
implementation and verified success rate.

\section{Reproducible artifacts and experimental boundary}
The exact density proof, PIC generator,
noise-model/count analyzer, incidence-transfer
audit, two-qutrit gate target and two-photon
bosonic postselection scripts are in
\texttt{analysis/w33\_*.py}. Their certificates
are frozen in \texttt{data/w33\_*.json};
related tests are in \texttt{tests/}.
The two FPGA compilation commands and
unpinned timing/resource ledger are in
\texttt{hardware/build\_w33\_walsh\_qutrit\_UP5K\_reference\_20261023.ps1}
and
\texttt{data/w33\_walsh\_qutrit\_UP5K\_timing\_20261023.json}.
The GDS's deliberately unresolved functions,
uncharacterized photonic noise, absent
experimental photon coincidence data and
lack of physical inter-bus optical control
are all explicit barriers. Nothing here
establishes a physical theory of everything
or a fault-tolerant quantum computer.
\end{document}
